\documentclass[12pt,a4paper]{ctexart}
\usepackage[margin=2.5cm]{geometry}
\usepackage{amsmath,amssymb,booktabs,graphicx}
\usepackage[hidelinks]{hyperref}
\usepackage[numbers,sort&compress]{natbib}

\title{数学建模论文}
\author{参赛队伍}
\date{}

\begin{document}
\maketitle

\begin{abstract}
本文档是可编译的中文数模论文骨架。正式写作时，正文和数值必须来自已经冻结并通过核验的写作包。

\textbf{关键词：}数学建模；可复现计算；鲁棒性
\end{abstract}

\section{问题重述}
根据赛题原文准确重述任务，不提前引入未经确认的方法假设。

\section{模型假设与符号}
列出经过确认的假设、符号及单位。

\section{模型建立与求解}
从验证后的分节文件引入最终方法、算法与结果。

\section{模型评价}
报告基线比较、敏感性、鲁棒性、局限和适用范围。

\section{结论}
仅陈述冻结结果支持的结论。

\end{document}
